\def\ChapterCode{GEF}
\chapter
[\CO{[GEF]} Gefahrenzonen, Gefahrentransporte, Gefahrengutunfall]
{Gefahrenzonen, \\Gefahrentransporte,
\\Gefahrengutunfall}
\label{topic:gefahrengutunfall}

\ChapterIntro
{%
~~~
}
{%
\begin{description}
  \item[Maintainer:] Roman Koch
  \item[Autoren:] Diverse
  \item[Reviewer:] Gerald Höritzmiller
  \item[Version:] {\VarDocVersionTypeName} ({\VarDocVersionTypeDescription})
  \item[SHA1:] {\DocGitCommit}
\end{description}

{\TxTeilkapitelAASS}
}


% \noindent\Parallel{
% % \includegraphics[width=\linewidth]{./images/noauth/AASdev20090228-img149.jpg}
% \includegraphics[width=\linewidth]{./icons/under-construction.png}
% }{
% % \includegraphics[width=\linewidth]{./images/noauth/AASdev20090228-img150.jpg}
% \includegraphics[width=\linewidth]{./icons/under-construction.png}
% }

\TODO{GABS}{2010-04-27}{Bilder 2x Gefahrengutunfall/LKW}{}

\clearpage

\section{Allgemeines}

Als Helfer steht man mitunter einer Vielzahl an möglichen
Gefahren gegenüber. Viele sind bereits durch überlegtes und
richtiges Handeln bewältigbar, andere bedürfen der
Unterstützung durch spezialisierte Kräfte. Kritisch ist das
rechtzeitige Erkennen von versteckten Gefahren, wie z.\,B.
Stick- und Giftgase, Chemikalien, usw.


\TextSlide{Persönliche Schutzausrüstung}
{\index{Schutzausrüstung!persönliche}%
Die \T{persönliche Schutzausrüstung} (\T{PSA}) Ist ein
wesentlicher Bestandteil zur Verhinderung von Verletzungen
und Erkrankungen des Personals im Einsatzgeschehen. Ihr
Umfang richtet sich nach den Anforderungen des jeweiligen
Einsatzes und den Umgebungsbedingungen. Zu der
\EE{Standardausrüstung} gehört eine den einschlägigen Normen
entsprechende \E{Arbeitsbekleidung} (lange Hose, \EE{Jacke
mit Ärmeln}), sowie \E{Sicherheitsschuhe} mit
durchtrittsicherer Sohle und Stahlkappen (Klasse S3).

Bei Patientenkontakt
werden vom Personal grundsätzlich \EE{Einweghandschuhe}
getragen. 

Die Arbeit auf
Verkehrsflächen erfordert eine Bekleidung mit
\EE{Warnschutz} (Tagesleuchtfarbe und Reflektorstreifen, gem.
\I{EN\,471}).\VARIANT{REG-AUT}{\footnote{In Österreich ist
das Tragen von Warnschutzbekleidung gesetzlich nur auf
Autobahnen festgeschrieben, dennoch ist das Tragen von
Warnschutzkleidung auch auf anderen Verkehrsflächen
sinnvoll.}} Für Tätigkeiten im Aussenbereich muss die
Kleidung über eine Wetterschutzfunktionalität gem. EN\,343
verfügen. Warn- und Wetterschutzbekleidungen sind mit
entsprechenden Symbolen gekennzeichnet, vgl.
\prettyref{fig:normen-warnschutz-wetterschutz}.


Bei Bedarf wird die Standardausrüstung durch weitere Ausrüstung
ergänzt, 
z.\,B. \EE{Helme} mit Nackenschutz,
 Einwegschürzen oder
\EE{Spritzschutzbrillen}. 
Bei infektiösen
Patienten werden je nach Erkrankung spezielle Anzüge bzw.
\EE{Atemmasken} verwendet. Als Atemmasken stehen
verschiedene Modelle zu Auswahl, von der
einfachen Mund-Nasen-Maske bis hin zu Partikelfiltermasken
der Klassen FFP~\VonBis{1}{3}, vgl.
\FullPrettyRef{topic:schutzkleidung-hygiene}, und
\FullPrettyRef{topic:atemmasken}.

In manchen Bereichen, z.\,B. im Wasser-, Berg-, Höhlen- oder
Pistenrettungsdienst, können die Erfordernisse erheblich vom
oben genannten abweichen.}
{
\begin{itemize}
  \item Standardausrüstung
  \begin{itemize}
    \item Arbeitskleidung: Lange Hose, Jacke mit Ärmeln
    \item Sicherheitsschuhe (S3)
    \item Einweghandschuhe
    \item Warnschutz
  \end{itemize}
  \item Weiters:
  \begin{itemize}
    \item Helme, Nackenschutz
    \item Einwegschürzen
    \item Spritzschutzbrillen
    \item Atemmasken
  \end{itemize}
  \item Weitere Ausstattung nach Bedarf
\end{itemize}
}
{}
{}
{}

\begin{Synopsis}
  \SynopsisItem{Die Standardschutzausrüstung, insbesonders die
  Arbeitsschutzjacke (inkl. der dazugehörigen Ärmel!), muss im
  Dienstbetrieb immer mitgeführt werden!}
\end{Synopsis}

\begin{PictureSeries}{}{Schutzausrüstung: Normen zum Warn-
und Wetterschutz\label{fig:normen-warnschutz-wetterschutz}}
\PictureSeriesRow{3}
{./images/pd/EN471}
{Kennzeichnung der Warnschutzeigenschaften nach EN\,471: X:
Fläche an fluroeszierndem Material (\VonBis{1}{3}); Y:
Fläche der Reflektorstreifen (\VonBis{1}{2})}
{./images/pd/EN343}
{Kennzeichnung der Schutzwirkung gegen Regen: X:
Wasserdichtigkeit (\VonBis{1}{3}); Y: Atmungsaktivität (\VonBis{1}{2})}
{empty}
{}
{}
{}
\end{PictureSeries}

\clearpage % TEMP
\section{Gefahrenzonen}
\label{topic:gefahrenzonen}

\TextSlide{Typische Gefahren}{%
Die Gefahren für den Helfer sind vielfältig. 
Allgemeinen, und damit \EE{alltäglichen}, Gefahrenquellen, 
wie z.\,B. dem Straßenverkehr,
wird in der Regel kompetent und routiniert begegnet.
%
Tückisch sind versteckte, bzw. nicht-alltägliche Gefahrenquellen, die sich jedoch unter Umständen fatal auswirken können.
Typische Beispiele sind 
\EE{Brand-} bzw. \EE{Explosionsgefahr},
\EE{giftige Stoffe}, 
\EE{Stickgase} (durch defekte Heizlüfter, Thermen; Gärgase, \ldots),
\EE{Infektionsgefahr} (Sonderformen der Grippe, hochinfektiöse Erkrankungen wie z.\,B. Ebola, Tubekulose, \ldots), 
das Entweichen unter \EE{Druck stehender Gase},
\EE{Ätzende Stoffe},
\EE{Radioaktivität},
\EE{Umweltgefährdung}.
%
Die Erfahrung zeigt, 
so selten manche dieser Situationen sein mögen, 
wenn man damit konfrontiert wird, 
geschieht dies meist unvorhergesehen.
%
Eine Gefahrensituation und Gefahrenzone rechtzeitig zu erkennen 
und frühzeitig die richtigen Maßnahmen zu ergreifen
ist wesentlich für die Bewältigung eines Einsatzes.
}{
\begin{itemize}
    \item Alltägliche Gefahren (Straßenverkehr, \ldots)
    \item Brand- bzw. Explosionsgefahr
    \item Infektionsgefahr
    \item Stickgase
    \item Giftige Stoff
    \item Entweichen unter Druck stehender Gase
    \item Ätzende Stoffe
    \item Radioaktivität
    \item Umweltgefährdung
\end{itemize}
}{}{}{}




\TextSlide{Typische Gefahrenquellen}{%
Auch die Entstehung von Gefahrenquellen ist vielfältig.
Besonders zu nennen sind Unfälle bei der \EE{Herstellung} von Stoffen (hier kann die Gefahr auch von Hilfsstoffen, z.\,B. Lösungsmitteln ausgehen), Unfällen bei der \EE{Verwendung} von gefährlichen Stoffen, sowie \EE{Transportunfälle}. 

Transportunfälle bergen besondere Gefahren, da sie sich normalerweise im ungesicherten Gelände ereignen, und die genaue Gefahrenlage anfangs schwer zu überblicken ist. Ebenso fehlt es, im Gegensatz zur chemischen Industrie, an fachkundigen Einweisern oder mit dem Gefahrgut geschultem Personal. 
}{
\begin{itemize}
    \item Unfälle bei der Herstellung
    \item Unfälle bei der Verwendung von gefährlichen
    Stoffen
    \item Transportunfälle
\end{itemize}
}{}{}{}


\TableWide
{Beispiele: Gefahrenquellen}%1 Titel
{}%2 Beschreibung
{}%3 Label
{}%4 Quelle
{}%5 Lizenz
{%
\begin{tabulary}{\linewidth}{LL}\toprule
\noindent\TabHeading{Situation}
&
\TabHeading{Typische Gefahren}
\\\midrule
\noindent\TabSub{Landwirtschaft}
 &
\ce{CO2} in Gärkellern und Weinkellern, Siloanlagen, Brunnenschächte

Pflanzenschutzmittel
\\\addlinespace
\noindent\TabSub{Klärgruben}
 &
\ce{CO2}, Methan, Schwefelwasserstoff
\\\addlinespace
\noindent\TabSub{Obstlagerhallen}
 &
wenig \ce{O2}, viel \ce{CO2} oder \ce{N2}, halten Obst frisch
\\\addlinespace
\noindent\TabSub{Brand}
 &
Rauchgase, extrem gefährlich besonders bei Wirtschafts- oder Industriegebäuden
\\\addlinespace
\noindent\TabSub{Tiefgaragen und Tunnels}
 &
erhöhte \ce{CO}- und \ce{CO2}-Werte
\\\addlinespace
\noindent\TabSub{Flüssiggas}
 &
Campingplätze, feste und mobile Küchen
\\\addlinespace
\noindent\TabSub{Pools, Schwimmbäder}
 &
Chlor (Wasserdesinfektion)
\\\addlinespace
\noindent\TabSub{Industrie}
 &
Kühlmittel in großen Kälteanlagen (Ammoniak), Lösungsmittel, gefährliche Hilfsstoffe
\\\addlinespace
\noindent\TabSub{Speziallabors, Einrichtungen mit und zur Verarbeitung von infektiösen
Abfällen oder Tierkadaver}
 &
Infektionsgefahr
\\\addlinespace
\noindent\TabSub{Nuklearmedizinische Institute, Industrie, Forschungseinrichtungen}
 &
Radioaktivität
\\\addlinespace
\noindent\TabSub{Transportunfälle}
 &
Brand- und Explosionsgefahr, Gift, Entweichung unter Druck stehender Gase, Infektionsgefahr, Radioaktivität, Umweltgefährdung, Straßenverkehr
\\\bottomrule
\end{tabulary}
}%6 Tabelle
{}%7 Float


\subsection{Elektrizität}

\Text{Strom -- allgemein}{%
\index{Gefahrenbereich!Strom}\label{topic:gefahrenbereich-strom}
%
Eine Starkstrom- bzw. Hochspannungsanlage muss durch einen
\EE{Fachmann} (E-Werk, ÖBB, Betriebsleitung, ...)
abgeschalten bzw. geerdet werden.

Die Anlage muss gegen Wiedereinschalten gesichert werden!
Patient mit isolierendem Material aus dem Stromkreis
befreien. Ein Sicherheitsabstand zu
stromführenden Leitungen muss eingehalten werden! 
}{%

}{}{}{}

\TextSlide{Freileitungen}{%
Bei Freileitungen mit
380\,000 Volt (380\,\MetricUnit{kV}) sind mindestens
20\Meter* Abstand zu halten.
\cite{BoehmerTaschRett6} Bei herabhängenden Hochspannungsleitungen oder
nach Blitzeinschlägen auf die Schrittspannung achten: Die
Spannung fällt im Boden nicht sofort auf 0~Volt ab, 
sondern bildet einen \T{Spannungskegel}. 
%
Innerhalb eines gewissen Umkreises kann daher, 
aufgrund des Potentialunterschiedes zwischen den beiden
Beinen am Boden, 
eine \T{Schrittspannung} auftreten.
%
Befindet man sich innerhalb eines solchen Gebietes, 
muss man die Beine eng beisammen halten 
und mit geschlossenen Beinen aus dem Gefahrenbereich herausspringen.
%
Je nach Bodenbeschaffenheit können Schrittspannungen
bis zu einer \EE{Distanz von 20\Meter*} von der herunterhängenden
Hochspannungsleitung auftreten. \cite{RedelsteinerHRNS1} 
%
Vor Betreten des Gefahrenbereiches muss der betreffende Energieversorger kontaktiert werden, 
und sichergestellt sein, 
dass die Leitung abgeschaltet wurde.
}{%
\begin{itemize}
  \item Spannungskegel
  \item Schrittspannung
  \item Abstand von 20\Meter* 
\end{itemize}
}{}{}{}

\subsection{Bahnanlagen}

\TextSlide{Bahnanlagen}{%

\begin{description}
  \item[ÖBB-Einsatzleiter] Der \EE{ÖBB-Einsatzleiter} ist
  der Ansprechpartner für die Sicherheit im Gleisbereich!
  Einsätze im Gleisbereich dürfen grundsätzlich nur mit
  Zustimmung des ÖBB-Einsatzleiters erfolgen!
  
  Bei \E{Gefahr in Verzug} dürfen alle erforderlichen
  Maßnahmen zur Menschenrettung nur unter Berücksichtigung
  der Eigensicherung (Einhaltung von Schutzabständen \ldots)
  und der Eigenverantwortung gesetzt werden.
  \cite{OeBBRettEinGleis2009}
  \item[Problem Zugverkehr] Bei
  Notfällen auf Bahngleisen ist auf jeden Fall der
  Streckenabschnitt von der ÖBB sperren zu lassen (über die
  Leitstelle)!
  \item[Problem Lokalisation] An den Masten der
  Oberleitungen der ÖBB sind Schilder mit Kilo- und
  Hektometerangaben angebracht\ZFootnote{Die obere Zahl gibt
  den Kilometer, die untere den Hektometer (100\Meter*)
  an.}. Diese müssen bei der Leitstelle gemeldet werden. Bei
  nicht-elektrifizierten Strecken finden sich die Angaben
  auf Hektometersteinen am Rand der Strecke. Weiters verfügt
  jeder Mast über eine Mastnummer, auch diese kann angegeben
  werden.
  \cite{OeBBRettEinGleis2009}
  
  \begin{Synopsis}
    \SynopsisItem{Die Nummernsysteme der ÖBB sind eine eigene
      Wissenschaft, im Zweifel sind zusätzlich zur Nummer auch
      anzugeben, woher diese stammt.}
  \end{Synopsis}
  
  \item[Problem Strom] Die
  Eisenbahn arbeitet mit einer Oberleitungsspannung von
  15\,000\,\MetricUnit{V}, auch bei gesenktem Stromabnehmer
  oder Dieselfahrzeugen kann durch Kondensatorladung eine
  Spannung von bis zu 3\,000\,\MetricUnit{V} anliegen.
  
  \VARIANT{VAR-STANDARD}{\Speech{\enquote{Die Freischaltung
  und Erdung wird vom ÖBB-Einsatzleiter veranlasst. Hilfs-
  und Einsatzkräfte erhalten beim ÖBB-Einsatzleiter vor
  Einsatzbeginn eine verbindliche Auskunft über den
  Schaltzustand der Oberleitung im Einsatzbereich. Alle
  Anlagenteile -- ausgenommen jene, die vom
  ÖBB-Einsatzleiter als freigeschaltet und geerdet
  bekanntgegeben werden -- sind als unter Spannung stehend
  zu betrachten!}}
  \cite{OeBBRettEinGleis2009}}
  
  Bei herabhängeenden Leitungen ist ein Sicherheitsabstand
  von mindestens 15\Meter* einzuhalten!
  
  Ähnliches gilt für Stadt- und U-Bahn-Anlagen. Da die
  Stromversorgung oft nicht durch eine Oberleitung, sondern
  durch andere Vorrichtungen in Bodennähe (Stromschiene,
  etc.) erfolgt, besteht auch bei intakten Anlagen die
  Gefahr eines Stromschlags, sofern die Gleisanlagen
  betreten werden. Als Traktionsspannung werden oft
  Spannungen um die 700\,\MetricUnit{V} eingesetzt.
  \cite{U61989}
  \item[Problem Bremsweg]
  Schienenfahrzeuge haben einen \EE{massiv verlängerten
  Bremsweg}: Ein Güterzug benötigt ca. 1\,\MetricUnit{km},
  ein Personenreisezug bis zu \EE{2,6\,\MetricUnit{km}}, um
  zum Stillstand zu kommen!
  
  \item[Problem Gefahrgut]
  Mit der Bahn transportiertes Gefahrgut wird analog zum
  Straßenverkehr mit Gefahrenzahl und UN-Nummer
  gekennzeichnet (\prettyref{topic:gefahrengutunfall}). Die
  zur Fracht gehörigen Frachtbriefe sind zu beachten.
\end{description}
}{%
\begin{itemize}
  \item ÖBB Einsatzleiter
  \item Zugverkehr
  \item Lokalisation: Hektometerschilder, -steine
  \item Strom:
  \begin{itemize}
    \item Bis zu 15\,000\,\MetricUnit{V}
    \item Erdung und Freigabe durch ÖBB-Einsatzleiter
  \end{itemize}
  \item Bremsweg: ↑↑↑
  \item Gefahrgut
\end{itemize}


}{}{}{}




\TextSlide{{Linktip}}{%
\begin{labeling}{mmm}
  \item[\cite{OeBBRettEinGleis2009}] \fullcite{OeBBRettEinGleis2009}
%   \emph{ÖBB Infrastruktur AG: Handbuch
%   Rettungseinsatz im ÖBB-Gleisbereich. online, Oktober 2009. \url{http://www.roteskreuz.at/uploads/media/Handbuch_Einsatz_Gleisbereich.pdf}}
\end{labeling}
}{\url{http://www.roteskreuz.at/uploads/media/Handbuch_Einsatz_Gleisbereich.pdf}
}{}{}{}















\section{Kennzeichnung von gefährlichen Stoffen und
Gefahrenbereichen}
\label{topic:gefahrsymbole}

\begin{PictureSeries}{}{Warnzeichen}
\PictureSeriesRow{4}
{./images/noauth/warnung-allgemein.png}
{Allgemeine Gefahr}
{./images/noauth/warnung-radioaktiv.png}
{Radioaktive Stoffe, ionisierende
Strahlen}
{./images/noauth/warnung-gasflaschen.png}
{Gasflaschen}
{./images/noauth/warnung-laser.png}
{Gebündeltes Licht / Laser}
\end{PictureSeries}


\subsection{Kennzeichnung von Stoffen}

\Text{Allgemeines}
{ %
Die Kennzeichnung von Gefahrenbereichen und gefährlichen Stoffen ist zwar normiert, 
aber abhängig vom Kontext 
(Kennzeichnung von Produktverpackungen bzw. Kennzeichnung von Gefahrguttransporten) 
und dem Zielpublikum 
(Feuerwehr, Endverbraucher, \ldots). 
%
Daher gibt es je nach Einsatzbereich unterschiedliche Kennzeichnungssysteme. 
}
{slide}

% \TableWide
% 	{<Titel>}%1
% 	{<Beschreibung>}%2
% 	{<Label>}%3
% 	{<Quelle>}%4
% 	{<Lizenz>}%5
% 	{<Tabelle>}%6
% 	{<Float>}%8
\TableWide
{Gefahrgutsymbole}%1
{Ab 1.12.2010}%2
{}%3
{Verordnung
(EG) Nr. 1272/2008 \cite{EG-1272-2008-De}}%4
{}%5
{
\FontSansNarrowFamily\scriptsize

\begin{tabularx}{\linewidth}{XXXXXXXXX}
\toprule\addlinespace
\multicolumn{9}{c}{{\normalsize\TabSub{Gefahrgutsymbole}}}
\\\addlinespace \multicolumn{9}{c}{\Z{\footnotesize nach Verordnung
(EG) Nr. 1272/2008 gültig ab 1.12.2010 \cite{EG-1272-2008-De}}}
\\\addlinespace\midrule\addlinespace

\includegraphics[width=\linewidth]{./images/GEF/gefahrensymbole/ghs01_explodingbomb.png}
&

\includegraphics[width=\linewidth]{./images/GEF/gefahrensymbole/ghs02_flame.png}
&

\includegraphics[width=\linewidth]{./images/GEF/gefahrensymbole/ghs03_flameovercircle.png}
&

\includegraphics[width=\linewidth]{./images/GEF/gefahrensymbole/ghs04_gascylinder.png}
&

\includegraphics[width=\linewidth]{./images/GEF/gefahrensymbole/ghs05_corrosion.png}
&

\includegraphics[width=\linewidth]{./images/GEF/gefahrensymbole/ghs06_skullandcrossbones.png}
&

\includegraphics[width=\linewidth]{./images/GEF/gefahrensymbole/ghs07_exclamationmark.png}
&

\includegraphics[width=\linewidth]{./images/GEF/gefahrensymbole/ghs08_healthhazard.png}
&

\includegraphics[width=\linewidth]{./images/GEF/gefahrensymbole/ghs09_environment.png}
\\\addlinespace
Gefahr:\newline
Explosions\-gefährlich
&
Gefahr:\newline
Leicht-/Hoch\-entzündlich 
&
Gefahr:\newline
Brandfördernd 
&
Achtung:\newline
Komprimierte Gase 
&
Gefahr:\newline
Ätzend 
&
Gefahr:\newline
Giftig/ Sehr giftig 
&
Achtung:\newline
Gesundheits\-gefährdend 
&
Gefahr:\newline
Gesundheits\-schädlich  
&
Warnung:\newline
Umwelt\-gefährdend
\\\addlinespace
\Z{\tiny GHS01}
&
\Z{\tiny GHS02}
&
\Z{\tiny GHS03}
&
\Z{\tiny GHS04}
&
\Z{\tiny GHS05}
&
\Z{\tiny GHS06}
&
\Z{\tiny GHS07}
&
\Z{\tiny GHS08}
&
\Z{\tiny GHS09}
\\\addlinespace\midrule\addlinespace
\multicolumn{9}{c}{Diese GHS-Symbole werden mit
Gefahrenhinweisen und Sicherheitshinweisen ergänzt!}
\\\addlinespace\toprule\addlinespace
\multicolumn{9}{c}{{\normalsize \TabSub{Alte Zeichen und
Gefahrenbezeichnungen} (Bis Mitte 2017)}}
\\\addlinespace\midrule\addlinespace
\includegraphics[width=\linewidth]{./images/GEF/gefahrensymbole/gefahr5.png}
&
\includegraphics[width=\linewidth]{./images/GEF/gefahrensymbole/gefahr1.png}
&
\includegraphics[width=\linewidth]{./images/GEF/gefahrensymbole/gefahr2.png}
&
&
\includegraphics[width=\linewidth]{./images/GEF/gefahrensymbole/gefahr3.png}
&
\includegraphics[width=\linewidth]{./images/GEF/gefahrensymbole/gefahr6.png}
&

&

&
\includegraphics[width=\linewidth]{./images/GEF/gefahrensymbole/gefahr4.png}
\\\addlinespace
\multicolumn{1}{c}{E} 
&
\multicolumn{1}{c}{F, F+} 
&
\multicolumn{1}{c}{O} 
&
&
\multicolumn{1}{c}{C} 
&
\multicolumn{1}{c}{T, T+} 
&

&

&
\multicolumn{1}{c}{N} 
\\\addlinespace
Explosions\-gefährlich 
&
Leicht-/Hoch\-entzündlich 
&
Brandfördernd 
&
keine Entsprechung 
&
Ätzend 
&
Giftig/ Sehr giftig
&
keine Entsprechung
&
keine Entsprechung 
&
Umwelt\-gefährlich
\\\addlinespace\bottomrule
\end{tabularx}
}%6




\subsection{Transport gefährlicher Güter}
\label{topic:adr}

\Text
{Einleitung}%1 Überschrift
{%
Fahrzeuge, die \E{ab einer bestimmten Menge} Güter
befördern, von denen eine Gefahr für Lebewesen und Umwelt
ausgehen kann (Gefahrgüter), müssen als solche deutlich
gekennzeichnet sein.
Die Kennzeichnung ermöglicht den Helfern bei einem Unfall
die richtigen Maßnahmen
treffen zu können. 
% (Klassifizierung nach ADR - Übereinkommen
% über die internationale Beförderung gefährlicher Güter auf
% der Straße).
Besondere Vorschriften für den Straßenverkehr hinsichtlich
Verpackung und Kennzeichnung von Gefahrgut sind im
\I[Europäisches Übereinkommen über die Beförderung
gefährlicher Güter auf der Straße]{Europäischen
Übereinkommen über die Beförderung gefährlicher Güter auf
der Straße} (\T{ADR}\ZFootnote{Accord européen relatif au transport
international des marchandises Dangereuses par Route} \citep{ADR2011})
geregelt.

Erst wenn die Menge eines transportierten Gefahrstoffes eine
festgelegte Mindestmenge überschreitet besteht eine
Kennzeichnungspflicht. Die Mindestmenge gilt jedoch
für jeden Stoff separat, und nicht für die Summe aller
transportierten Stoffe.
Es können also erhebliche Mengen gefährlicher Güter ohne
Kennzeichnung des Fahrzeuges transportiert werden. 
}%2 Text





\Text{Gefahrentafel}
{\label{topic:warntafel}
Die Gefahrentafel ist eine 40\,$\times$\,30\CM* große
orangefarbene Tafel mit schwarzem Rand, welche auf dem
Fahrzeug angebracht ist und einen Gefahrguttransport
kennzeichnet. 
An ihr kann eine \I{Gefahrnummer} (auch: \I{Kemler-Nummer})
obere Hälfte) und eine Stoffnummer (auch: \I{UN-Nummer},
untere Hälfte) angebracht sein.
}
{}
{}
{}
{}



\TableWide
{Kennzeichnung durch Gefahrentafel}%1 Titel
{}%2 Beschreibung
{}%3 Label/Hook
{}%4 Quelle
{}%5 Lizenz
{%
\begin{tabular}{m{3.225cm}m{3.225cm}m{3.225cm}m{3.225cm}}
\toprule\addlinespace
\multicolumn{2}{m{6.65cm}}{\centering \TabHeading{Allgemeine Kennzeichnung}
(Sammeltransporte)} &
\multicolumn{2}{m{6.65cm}}{\centering \TabHeading{Spezielle Kennzeichnung}} 
\\\addlinespace\midrule\addlinespace
\includegraphics[width=\linewidth]{./images/khd/warntafel-allgemein.png}
 &
 &
\includegraphics[width=\linewidth]{./images/khd/warntafel-un-60-1710.png}
 &
Gefahrnummer\par (\T{Kemler-Nummer})

Stoffnummer\par (\T{UN-Nummer})
\\\addlinespace\bottomrule
\end{tabular}
}%6 Tabelle
{}%7 Float
% vormals \Layout{57}

\Text{Gefahrnummer}
{\label{topic:gefahrnummer}
Die Gefahrnummer gibt Aufschluss über die zu erwartenden
Gefahren eines transportierten Stoffes und ist im oberen
Teil der Gefahrentafel angebracht. Sie setzt sich aus
Ziffern zusammen, welche jeweils eine spezielle Bedeutung
haben (\FullPrettyRef{tab:gefahrennummer}). Die Verdopplung
einer Ziffer weist auf die Zunahme der entsprechenden Gefahr
hin. Wenn die Gefahr eines Stoffes ausreichend von einer
einzigen Ziffer angegeben werden kann, wird dieser Ziffer
eine Null angefügt. 

\begin{Danger}
  \item Gefahrnummer: Ein vorangestelltes \enquote*{X} bedeutet,
  dass der Stoff in  gefährlicher Weise
  \Ca{mit Wasser reagiert}!
  
  \Ca{Vorsicht: Nebel, Regen, Schnee, Löschwasser}
\end{Danger}

% 
% Die in
% \FullPrettyRef{tab:gefahrennummer-spezielle-ziffernkombinationen},
% aufgeführten Ziffernkombinationen haben eine besondere
% Bedeutung.

}
{}
{}
{}
{}



\Text{Stoffnummer}
{Die Stoffnummer ist ein Code, welcher Aufschluss über den
transportierten Stoff gibt. Kataloge zur Entschlüsselung
liegen i.\,d.\,R. in der Leitstelle und den Fahrzeugen der
Feuerwehr auf.

}
{}
{}
{}
{}



\Table
{Bedeutung der
Ziffern der Gefahrennummer}%1 Titel
{}%2 Beschreibung
{\label{tab:gefahrennummer}}%3 Label
{}%4 Quelle
{}%5 Lizenz
{%
\begin{tabulary}{\linewidth}{cL}
\addlinespace\toprule\addlinespace
 2 & Entweichen von Gas durch Druck oder chemische Reaktion
\\\addlinespace
 3 & Entzündbarkeit von Flüssigkeiten (Dämpfen) und Gasen oder 
selbsterhitzungsfähiger flüssige
Stoffe (z.B. Benzin, Diesel)
\\\addlinespace
 4 & Entzündbarkeit von festen Stoffen oder
selbsterhitzungsfähiger fester Stoffen
\\\addlinespace
 5 & Oxidierende (brandfördernde) Wirkung
\\\addlinespace
 6 & Giftigkeit oder Ansteckungsgefahr
\\\addlinespace
 7 & Radioaktivität
\\\addlinespace
 8 & \"Atzwirkung
\\\addlinespace
 9 & Gefahr einer spontan heftigen Reaktion
\\\addlinespace
0 & Keine weitere Gefahr
\\\addlinespace
 {\sffamily X} & Reagiert gefährlich mit Wasser
\\\addlinespace\bottomrule
\end{tabulary}
}%6 Tabelle
{}%7 Float



\Table
{Spezielle Ziffernkombinationen bei der
Gefahrennummer}%1 Titel
{}%2 Beschreibung
{\label{tab:gefahrennummer-spezielle-ziffernkombinationen}}%3 Label
{}%4 Quelle
{}%5 Lizenz
{%
\begin{tabulary}{\linewidth}{cL}
\addlinespace\toprule\addlinespace
22 & Tiefgekühltes Gas
\\\addlinespace
X323 & Entzündbarer flüssiger Stoff, der mit Wasser gefährlich reagiert, wobei entzündbare Gase entweichen
\\\addlinespace
X333 & Selbstentzündliche Flüssigkeit, die mit Wasser gefährlich reagiert
\\\addlinespace
X423 & Entzündbarer fester Stoff, der mit Wasser gefährlich reagiert, wobei brennbare Gase entweichen
\\\addlinespace
44 & Entzündbarer fester Stoff, der sich bei erhöhter Temperatur in Geschmolzenem Zustand befindet
\\\addlinespace
539 & Entzündbares organisches Peroxid
\\\addlinespace
90  & Verschiedene gefährlich Stoffe
\\\addlinespace\bottomrule
\end{tabulary}
}%6 Tabelle
{}%7 Float







\Table
{Ausgewählte Stoffnumern}%1 Titel
{}%2 Beschreibung
{}%3 Label/Hook
{}%4 Quelle
{}%5 Lizenz
{%
\begin{tabulary}{\linewidth}{cLcL}
\toprule\addlinespace
0048 & Sprengkörper & 1223 & Kerosin
\\\addlinespace
0333 & Feuerwerkskörper & 1789 & Salzsäure
\\\addlinespace
1005 & Ammoniak & 1830 & Schwefelsäure
\\\addlinespace
1011 & Butan & 1962 & Ethylen
\\\addlinespace
1013 & Kohlendioxid & 1972 & Methan oder Erdgas (tiefgekühlt, flüssig)
\\\addlinespace
1017 & Chlor & 1977 & Stickstoff
\\\addlinespace
1052 & Flußsäure & 1978 & Propan
\\\addlinespace
1090 & Aceton & 3295 & Kohlenwasserstoff, flüssig
\\\addlinespace
1203 & Benzin & 3312 & Gas, tiefgekühlt, flüssig
\\\addlinespace\bottomrule
\end{tabulary}
}%6 Tabelle
{}%7 Float






\Text{Gefahrzettel}
{\label{topic:gefahrzettel}
Zusätzlich zur Kennzeichnung mittels der Gefahrennummer 
erfolgt eine Kennzeichnung mittels Gefahrzettel.
Dies sind farbige Piktogramme in Form eines auf der Spitze stehenden Quadrates. 
%
Im oberen Bereich ist ein die Gefahr anzeigendes Symbol angebracht,
in der unteren Hälfte findet sich die Nummer der Gefahrenklasse.
%
Bei Stückguttransporten wird jedes Versandstück
gekennzeichnet. 
%
An Fahrzeugen, die kein Stückgut
transportieren, müssen Gefahrenzettel an den Längsseiten und
an der Rückseite angebracht werden.}
{}
{}
{}
{}



\TableWide
{Gefahrenzettel}%1 Titel
{\EE{Gefahrenklassen} und ausgewählte Symbole nach \citep{ADR2011}}%2 Beschreibung
{}%3 Label/Hook
{}%4 Quelle
{}%5 Lizenz
{%
\begin{tabular}{m{0.3\linewidth}m{0.15\linewidth}m{0.3\linewidth}m{0.15\linewidth}}
\toprule\addlinespace
% \TabHeading{Gefahrenklasse} 
% &
% \TabHeading{Symbol}
% 
% (Auswahl)
% \\\addlinespace\midrule\addlinespace
\TabSub{Klasse 1}

Explosive Stoffe und Gegenstände mit Explosivstoff

Unterklassen:\newline
1.1, 1.2, 1.3, 1.4, 1.5, 1.6 
&
~\newline
\includegraphics[width=1.5cm]{./images/GEF/gefahrgut/gefahrgut_1.png}
\\\addlinespace\midrule\addlinespace
\TabSub{Klasse~2.1}

Entzündbare Gase 
&
~\newline
\includegraphics[width=1.5cm]{./images/GEF/gefahrgut/gefahrgut_21.png}
&
\TabSub{Klasse~2.2}

Nicht entzündbare,\newline
nicht giftige Gase 
&
~\newline
\includegraphics[width=1.5cm]{./images/GEF/gefahrgut/gefahrgut_22.png}
\\
\TabSub{Klasse~2.3}

Giftige Gase 
&
~\newline
\includegraphics[width=1.5cm]{./images/GEF/gefahrgut/gefahrgut_23.png}
 \\\addlinespace\midrule\addlinespace
\TabSub{Klasse~3}

Entzündbare\newline
flüssige Stoffe 
&
~\newline
\includegraphics[width=1.5cm]{./images/GEF/gefahrgut/gefahrgut_3.png}
\\\addlinespace\midrule\addlinespace
\TabSub{Klasse~4.1}

Entzündbare feste Stoffe, selbstzersetzliche Stoffe und
desensibilisierte explosive feste Stoffe 
&
~\newline
\includegraphics[width=1.5cm]{./images/GEF/gefahrgut/gefahrgut_41.png}
&
\TabSub{Klasse~4.2}

Selbstentzündliche Stoffe 
&
~\newline
\includegraphics[width=1.5cm]{./images/GEF/gefahrgut/gefahrgut_42.png}
\\
\TabSub{Klasse~4.3}

Stoffe, die in Berührung mit Wasser entzündliche Gase entwickeln 
&
~\newline
\includegraphics[width=1.5cm]{./images/GEF/gefahrgut/gefahrgut_43.png}
\\\addlinespace\midrule\addlinespace
\TabSub{Klasse~5.1}

Entzündend\newline
(oxidierend)\newline
wirkende Stoffe 
&
\includegraphics[width=1.5cm]{./images/GEF/gefahrgut/gefahrgut_51.png}
&
\TabSub{Klasse~5.2}

Organische Peroxide
&
~\newline
\includegraphics[width=1.5cm]{./images/GEF/gefahrgut/gefahrgut_52.png}
\\\addlinespace\midrule\addlinespace
\TabSub{Klasse~6.1}

Giftige Stoffe 
&
~\newline
\includegraphics[width=1.5cm]{./images/GEF/gefahrgut/gefahrgut_61.png}
&
\TabSub{Klasse~6.2}

Ansteckungs\-gefährliche Stoffe 
&
~\newline
\includegraphics[width=1.5cm]{./images/GEF/gefahrgut/gefahrgut_62.png}
\\\addlinespace\midrule\addlinespace
\TabSub{Klasse~7}

Radioaktive Stoffe 
&
\includegraphics[width=1.5cm]{./images/GEF/gefahrgut/gefahrgut_7.png}
~
\includegraphics[width=1.5cm]{./images/GEF/gefahrgut/gefahrgut_7_fissile.png}
\\\addlinespace\midrule\addlinespace
\TabSub{Klasse~8}

Ätzende Stoffe 
&
\includegraphics[width=1.5cm]{./images/GEF/gefahrgut/gefahrgut_8.png}
\\\addlinespace\midrule\addlinespace
\TabSub{Klasse~9}

Verschiedene gefährliche Stoffe und Gegenstände &


\includegraphics[width=1.5cm]{./images/GEF/gefahrgut/gefahrgut_9.png}
\\\addlinespace\midrule\addlinespace
\end{tabular}
}%6 Tabelle
{}%7 Float

% \Layout{57}{}{}{}{
% \begin{tabular}{m{2.cm}m{3.5cm}m{2.cm}m{3.5cm}}
%  Beispiele
%  &
%  &
%  &
% \\
% \includegraphics[width=1.659cm,height=1.885cm]{./images/noauth/gefahrzettel-feuergefaehrlich-fluessig.png}
%  &
% Feuergefährlich \par(flüssige Stoffe)
%  &
% \includegraphics[width=1.856cm,height=2.096cm]{./images/noauth/gefahrzettel-selbstentzuendlich.png}
%  &
% Selbstentzündlich
% \\
% \includegraphics[width=1.835cm,height=1.902cm]{./images/noauth/gefahrzettel-feuergefaehrlich-fest.png}
%  &
% Feuergefährlich \par(feste Stoffe)
%  &
% \includegraphics[width=1.909cm,height=2.162cm]{./images/noauth/gefahrzettel-enzuendlich-wasser.png}
%  &
% Entzündbare Gase bei Berührung mit \EE{Wasser}
% \\
% \end{tabular}
% }{}{}












\clearpage % TEMP temp clearpage
\section{Allgemeine Einsatzrichtlinie}
\label{topic:einsatzrichtlinie}


\TextSlide{Übersicht - GAS-Regel}{

\begin{itemize}
  \item \EE{G}efahr erkennen
  \item \EE{A}bsperrmaßnahmen durchführen
  \item \EE{S}pezialkräfte anfordern
\end{itemize}
}{
\begin{itemize}
  \item \EE{G}efahr erkennen
  \item \EE{A}bsperrmaßnahmen durchführen
  \item \EE{S}pezialkräfte anfordern
\end{itemize}
}{}{}{}

\begin{PictureSeriesWide}{}{GAS-Regel
\SourcePicture{Gabriel}}
\PictureSeriesRowWide{3}
{./images/gabriel-sebastian-aass/UE2011FORTUNA-00268-0154pt-0300dpi}
{\TabHeading{{\color{ubuntured}G}efahr erkennen}: Wenn in
einem Raum alle Personen bewusstlos sind, dann ist das
eigenartig \ldots }
{./images/gabriel-sebastian-aass/UE2011FORTUNA-00964-0154pt-0300dpi}
{\TabHeading{{\color{ubuntured}A}bsperr\-maßnahmen durch\-führen}
}
{./images/gabriel-sebastian-aass/UE2011FORTUNA-00518-0154pt-0300dpi}
{\TabHeading{{\color{ubuntured}S}pezial\-kräfte anfordern}:
Informationen angeben (Kemmler-Nummer, etc.)
}
{}
{}
\end{PictureSeriesWide}




\Text{Gefahr erkennen}
{Das Erkennen der Gefahr ist Grundvoraussetzung um richtig
zu handeln. Im Idealfall wird eine Gefahr beim
Szeneüberblick erkannt, manchmal jedoch erst während der
Patientenversorgung.
Mögliche Hinweise sind z.\,B.:
\begin{itemize}
  \item Einsatzort 
  \item Situation
  \item Hinweisschilder
  \item Kennzeichnung mittels Gefahrentafel, Gefahrnummer
  (Kemler-Nummer), Stoffnummer (UN-Nummer), Gefahrensymbole,
  allgemeiner Warnzeichen
  \item \IX{Frachtbriefe}
  \item \ldots
\end{itemize}

}{}
{}
{}
{}


\Text{Absperrmaßnahmen durchführen}
{\label{topic:absperrmassnahmen}
Die Absperrung sollte großräumig  sein, genügend Platz für
die technischen Hilfeleistungen (Feuerwehr) muss
einkalkuliert werden. Die Witterung, insbesonders
Windrichtung und Niederschlag, muss bedacht werden. Es muss
mit sich ändernden Situationen gerechnet werden
(Schaulustige, Änderung der Windrichtung, Böen, \ldots).

}
{}
{}
{}
{}

\begin{PictureSeriesWide}{}
{Bilderserie:
\HeadingKeyword{Misslungene Absperrung} \SourcePicture{Gabriel}}
\PictureSeriesRowWide{3}
{./images/khd/gef-wind-rauch1.jpg}
{Ein Auto brennt und alle schauen
fasziniert zu \ldots }
{./images/khd/gef-wind-rauch2.jpg}
{\ldots doch dann kommt Wind auf \ldots  }
{./images/khd/gef-wind-rauch.jpg}
{Und oh Schreck! Die Absperrung
verliert ihren Zweck! }
{}
{}
\end{PictureSeriesWide}

\begin{Synopsis}
  \SynopsisItem{Die Moral aus der Geschicht': Vertrau' dem lauen Klima nicht!}
\end{Synopsis}



\Text{Spezialkräfte anfordern}
{Bei Gefahrgut und -stoffen ist häufig Hilfe von
Spezialisten gefragt. In der Regel ist hier die Feuerwehr
zuständig, bzw. kann weitere Hilfe vermitteln.
}
{}
{}
{}
{}




\nocite{Cicero911CommissionReport}



\ChapterEnd
